% =============================================================================
% style-module-architecture.tex
% Style: ltechnical-design-spec (which transitively loads
%        latex-docs-technical -> base.sty -> latex-docs-style).
% Build: pdflatex via latexmk; minted-aware with listings fallback.
% =============================================================================
\documentclass[11pt]{article}
\usepackage{ltechnical-design-spec}

\setDocTitle{Modular LaTeX style-module architecture}
\setDocSubtitle{Inheritance-oriented design for the \texttt{latex-docs} repository}
\setDocVersion{v1.0}
\setDocStatus{Accepted}
\setDocOwner{Tooling \& Documentation}
\setDocAuthor{Jordan Suber}
\setDocSummary{%
  A reusable, hierarchical LaTeX style system organised by document
  function. A single parent module owns shared presentation foundations;
  family modules (legal, financial, business / administrative, HR,
  personal / official, and an abstract technical) generalise from the
  parent; five technical subtypes generalise from the abstract technical
  module. The system preserves backward compatibility with the existing
  canonical house style and is designed for long-term reuse.%
}

\usepackage{tikz}
\usetikzlibrary{positioning,arrows.meta,shapes.geometric,fit,backgrounds}

\begin{document}
\makedoctocpage

% =============================================================================
\section{Proposed style-module hierarchy}

The hierarchy is a single-rooted tree with one abstract intermediate node
(the technical family). Every concrete document module is a leaf; every
non-leaf is either the parent or the abstract intermediate. Backward
compatibility with the canonical house style is preserved at the parent
level by delegation.

\begin{center}
\resizebox{\textwidth}{!}{%
\begin{tikzpicture}[
  node distance=4mm and 8mm,
  every node/.style={font=\small},
  abs/.style    ={draw=Steel, dashed, fill=Soft,    rounded corners=2pt,
                  inner sep=4pt, minimum width=44mm, align=center},
  fam/.style    ={draw=Navy,           fill=Soft,    rounded corners=2pt,
                  inner sep=4pt, minimum width=44mm, align=center},
  sub/.style    ={draw=Teal,           fill=white,   rounded corners=2pt,
                  inner sep=3pt, minimum width=58mm, align=center},
  tech/.style   ={draw=Steel, dashed, fill=Soft!80, rounded corners=2pt,
                  inner sep=4pt, minimum width=44mm, align=center},
  edge/.style   ={-{Latex[length=2mm]}, Steel, line width=0.5pt}
]
% Parent
\node[abs]                                                       (root) {%
  \textbf{base.sty}\\\scriptsize\itshape parent (abstract)};

% Family row
\node[fam, below=12mm of root, xshift=-72mm] (legal)            {legal};
\node[fam, right=2mm of legal]              (fin)               {financial};
\node[fam, right=2mm of fin]                (biz)               {business-admin};
\node[fam, below=4mm of legal, xshift=24mm]  (hr)                {hr};
\node[fam, right=2mm of hr]                 (po)                {personal-official};
\node[tech, right=2mm of po]                (tech)              {%
  \textbf{technical}\\\scriptsize\itshape abstract intermediate};

% Edges parent -> families
\foreach \n in {legal, fin, biz, hr, po, tech}
  \draw[edge] (root.south) -- ($(\n.north)+(0,1pt)$);

% Subtype row
\node[sub, below=14mm of tech, xshift=-46mm]      (um)  {technical-user-manual};
\node[sub, right=1mm of um]                      (ds)   {technical-design-spec};
\node[sub, below=2mm of um, xshift=22mm]         (impl) {technical-implementation};
\node[sub, below=2mm of impl, xshift=-30mm]      (inst) {technical-installation};
\node[sub, right=1mm of inst]                    (test) {technical-testing};

% Edges technical -> subtypes
\foreach \n in {um, ds, impl, inst, test}
  \draw[edge] (tech.south) -- (\n.north);

\end{tikzpicture}%
}
\end{center}

\Constraint{Inheritance is single, semantic, and one-way:} children
generalise from one parent. Documents never load \Module{base.sty}
directly and never load two sibling modules simultaneously.

% =============================================================================
\section{Parent module responsibilities}

\Module{base.sty.sty} is the abstract root of the hierarchy.
Concretely it is responsible for:

\begin{itemize}
  \item \textbf{Foundation delegation.} It \texttt{\textbackslash RequirePackage}s the
        canonical house style \Module{latex-docs-style.sty} once. Page
        geometry, font selection, microtypography, the colour palette
        (Navy / Steel / Teal / Brick on Soft), tables / lists / floats,
        the standard tcolorbox callouts (note / tip / warn / info / view),
        hyperlinks, code rendering, metadata API, page style, and
        section-heading typography all come from there. Family modules
        inherit them transitively, never re-declare them.
  \item \textbf{Module-classification registry.}
        \Function{LdsRegisterModule} sets two macros every consumer can
        read: \Module{LdsCurrentModuleKey} and \Module{LdsCurrentModuleName}.
        Subtypes call it after their parent's call, so the deepest module
        wins.
  \item \textbf{\Module{DocClass} metadata.} A new metadata field carrying
        the document's functional class. It flows automatically into the
        title-block meta line and the running header via
        \Function{LdsExtraHeader}.
  \item \textbf{Title and header extension hooks.}
        \Function{LdsExtraTitleMeta} and \Function{LdsExtraHeader} let
        children augment the title block and header without
        re-implementing them.
  \item \textbf{Generic counter factory.} \Function{NewLdsCounter} declares
        a counter, configures \Module{cleveref} singular / plural names,
        and registers a \Function{LdsRefCounter} accessor. Every family
        uses it; no family redefines counter glue from scratch.
  \item \textbf{Labeled-block environment factory.}
        \Function{NewLdsLabeledBlock} stamps out a tcolorbox-backed,
        counter-driven, optionally-titled environment in a single line.
        Three legal callouts, four installation steps, an ADR header,
        and the test-case header all use the same factory.
  \item \textbf{Shared signature primitive.} \Function{LdsSignatureLine}
        is composed by legal, HR, and personal-official modules.
        Composition over re-implementation.
  \item \textbf{Hairline-overfull tolerance.} A modest
        \Module{emergencystretch} so headings combining bold + small caps +
        colour absorb tight measure without warning.
\end{itemize}

\begin{notebox}[Inheritance vs. Composition]
The shared signature primitive deserves a note. Three families need a
signature line: legal, HR (offer letters, performance reviews),
and personal-official (certificates). They are not a single inheritance
chain; the right primitive is therefore \emph{composed in} from the
parent and re-used. Inheritance models specialisation; composition
models a shared-but-orthogonal capability. The architecture uses both.
\end{notebox}

% =============================================================================
\section{Function-class module responsibilities}

Each top-level family module owns the structures that ALL documents in
its functional class use, and nothing more. The litmus test for adding
behaviour to a family module: ``Will every document of this kind want
this, or only some?'' If only some, the behaviour belongs in the
document, not the module.

\begin{traceability}
  \Trace{Legal}{Articles, clauses, recitals, parties, definitions, signatures.}{latex-docs-legal}{example-legal-contract}
  \Trace{Financial}{Currency, line items, summary rows, ledger, invoice header.}{latex-docs-financial}{example-financial-invoice}
  \Trace{Business / Admin}{Memo header, attendees, agenda, decisions, action items, policies, procedures.}{latex-docs-business-admin}{example-business-memo}
  \Trace{HR}{Job descriptions, applicant records, training schedules, performance rubrics, PII redaction.}{latex-docs-hr}{example-hr-job-description}
  \Trace{Personal / Official}{Certificate page, ID-card layout, single-signature primitive.}{latex-docs-personal-official}{example-personal-certificate}
  \Trace{Technical (abstract)}{Requirements, long commands, hazard / caution callouts, step-result tables, semantic tagging.}{latex-docs-technical}{(intermediate; loaded transitively)}
\end{traceability}

% =============================================================================
\section{Technical-document specialization model}

The technical family is the only one with internal structure. Modelling
the five technical subtypes as siblings under a shared abstract module
captures the empirical reality that user manuals, design specs,
implementation guides, installation runbooks, and test plans share a
substantial substrate -- requirement IDs, command rendering, hazard
callouts, step-and-expected-outcome tables -- while each adds genuinely
distinct structures.

\begin{adr}{Make \texttt{latex-docs-technical} an abstract intermediate}{Accepted}
  \ADRContext{%
    Five technical document subtypes need shared substrate but each adds
    distinct structures. Two designs were considered: (a) flatten all
    five as direct children of the parent, with shared substrate
    duplicated in each; (b) introduce an abstract intermediate technical
    module that owns the substrate, with concrete subtypes generalising
    from it.%
  }
  \ADRDecision{%
    Adopt (b). The abstract \Module{latex-docs-technical} is loadable on
    its own (so a generic technical document still has a home) but the
    recommended endpoint for a real document is one of the five concrete
    subtypes. The intermediate is the natural place for any future
    addition (e.g.\ release notes, architecture overview) that fits the
    technical family but does not match an existing subtype.%
  }
  \ADRConsequences{%
    The hierarchy gains one node, a small price for substantial
    code reuse. Subtype implementation is a few dozen lines instead of
    several hundred. Adding a new technical subtype is a 30-minute task,
    not a 4-hour task.%
  }
\end{adr}

\subsection{What each subtype contributes}

\begin{description}
  \item[\Module{technical-user-manual}] \texttt{task} and \texttt{tutorial}
    environments, GUI semantics (\Function{Menu}, \Function{Button},
    \Function{Field}, \Function{Tab}, \Function{Window}),
    \texttt{troubleshooting} Q/A table, reading-level-tuned note wrappers.
  \item[\Module{technical-design-spec}] \Function{ADR} header with status
    pill, \texttt{adr} environment with Context / Decision / Consequences,
    \texttt{traceability} matrix, numbered \Function{DesignDecision}
    callouts.
  \item[\Module{technical-implementation}] \texttt{apicontract} environment
    with Inputs / Outputs / Errors / Side effects, \Function{CodeRef}
    inline source pointers, \texttt{walkthrough} narrative environment,
    semantic source-element typography (\Function{Module},
    \Function{Function}, \Function{TypeName}).
  \item[\Module{technical-installation}] \texttt{prerequisites} checklist,
    \Function{IStep} step blocks, \Function{Verify} and
    \Function{RollbackNote} step-local secondary fields, top-level
    \texttt{rollback} callout, \Function{EnvVar} typography.
  \item[\Module{technical-testing}] \Function{TestCase} header,
    \texttt{testcase} environment, status glyphs (\Function{Pass},
    \Function{Fail}, \Function{Skip}, \Function{Block}),
    \texttt{teststeps} table, \texttt{coverage} requirements-to-tests
    matrix.
\end{description}

% =============================================================================
\section{Recommended file / folder structure}

\begin{minted}[fontsize=\footnotesize]{text}
tooling/latex/
  latex-docs-style.sty              # canonical house style (UNCHANGED)
  styles/
    README.md                       # quick reference
    base.sty.sty             # parent / hierarchy root
    legal.sty
    financial.sty
    business-admin.sty
    hr.sty
    personal-official.sty
    technical.sty                   # abstract intermediate
    technical-user-manual.sty
    technical-design-spec.sty
    technical-implementation.sty
    technical-installation.sty
    technical-implementation.sty

src/architecture/style-system/
  style-module-architecture.tex     # this document
  examples/
    example-legal-contract.tex
    example-financial-invoice.tex
    example-business-memo.tex
    example-hr-job-description.tex
    example-personal-certificate.tex
    example-technical-user-manual.tex
    example-technical-design-spec.tex
    example-technical-implementation.tex
    example-technical-installation.tex
    example-technical-testing.tex
\end{minted}

\subsection{Naming convention}

Every module file is \Module{latex-docs-<segment>[-<segment>]*.sty} so all
modules sort together in \cmd{kpsewhich -all}, in IDE autocomplete, and in
the \texttt{texmf} tree. Module identifiers (the argument passed to
\Function{LdsRegisterModule}) are kebab-case and match the file stem
without the prefix: \texttt{technical-user-manual},
\texttt{technical-installation}, etc.

The repo's \texttt{Makefile} exports
\Cmdline{TEXINPUTS=\$(abspath tooling/latex)//:\$(TEXINPUTS)}
where the trailing \texttt{//:} makes \texttt{kpathsea} search recursively.
Files under \texttt{tooling/latex/styles/} are therefore discovered
automatically; no \texttt{Makefile} change is needed when new modules are
added.

% =============================================================================
\section{Sample \texttt{.sty} skeletons}

The canonical skeletons are the actual modules in
\Cmdline{tooling/latex/styles/}. Reading those is the authoritative
source. The shape of any new family module is:

\begin{minted}[fontsize=\footnotesize]{latex}
\NeedsTeXFormat{LaTeX2e}
\ProvidesPackage{latex-docs-<family>}%
  [<YYYY/MM/DD> v<n> latex-docs <family> family module]
\RequirePackage{base.sty}
\LdsRegisterModule{<family>}{<Display Name>}

% Counters via the parent factory; no manual cleveref / refstepcounter glue.
\NewLdsCounter{<Short>}{<crefnameplural>}{<crefname>}

% Labeled-block environments via the parent factory.
\NewLdsLabeledBlock{<envname>}{<counterShort>}{<word>}{<frame colour>}

% Family-specific environments and macros.
\newenvironment{<env>}{...}{...}
\newcommand{\<Macro>}[<n>]{...}

\endinput
\end{minted}

The shape of a technical subtype is identical except its first non-comment
line is \Cmdline{\textbackslash RequirePackage\{latex-docs-technical\}}
instead of \Cmdline{\textbackslash RequirePackage\{base.sty\}}.

% =============================================================================
\section{Sample \texttt{.tex} usage examples}

Ten examples ship in \Cmdline{src/architecture/style-system/examples/}, one
per concrete module. Each is a complete, compilable, single-file document.
The canonical pattern is:

\begin{minted}[fontsize=\footnotesize]{latex}
\documentclass[11pt]{article}
\usepackage{latex-docs-<family-or-subtype>}

\setDocTitle{...}        \setDocSubtitle{...}
\setDocVersion{...}      \setDocStatus{...}
\setDocOwner{...}        \setDocAuthor{...}
\setDocDate{\today}      \setDocSummary{...}

\begin{document}
\makedoctocpage  % or \makedoctitle alone for short documents
% ...content using the family's specialised environments...
\end{document}
\end{minted}

\Function{DocClass} is set automatically by the loaded module; documents
need not set it. To override (e.g.\ to refine a subtype label), call
\Cmdline{\textbackslash setDocClass\{Technical / Installation -- Cluster v3\}}
in the preamble.

% =============================================================================
\section{Design rationale and extension guidance}

\subsection{What belongs at each level}

\begin{description}
  \item[Parent] Anything genuinely shared by ALL function families.
    The litmus test: ``If we removed family X, would the parent still
    need this?'' If yes, it belongs in the parent. If no, it does not.
  \item[Family module] Anything genuinely shared by ALL documents of a
    function class -- the structures that define the class. If only a
    subset of documents in the class uses it, it does not belong here;
    push it to a subtype or to the document.
  \item[Subtype] Structures unique to the subtype. ADR shape lives only
    in design-spec. Test-case status glyphs live only in testing. The
    subtype is also the right home for narrowly-scoped typographic
    conventions (\Function{Menu} / \Function{Button} only make sense
    in user manuals, not in test plans).
\end{description}

\subsection{Avoiding duplication and fragile inheritance}

\begin{itemize}
  \item \textbf{Use factories, not copy-paste.} \Function{NewLdsCounter} and
        \Function{NewLdsLabeledBlock} are deliberately generic so families
        do not re-implement counter / cleveref / tcolorbox glue.
  \item \textbf{Use composition for cross-family capability.}
        \Function{LdsSignatureLine} is composed in by legal, HR, and
        personal-official because signatures are not a single inheritance
        chain. Inheritance models specialisation; composition models a
        shared-but-orthogonal capability.
  \item \textbf{Never override foundations.} A child module never
        \Function{renewcommand}s a foundation macro from
        \Module{latex-docs-style}. Children extend through the hooks
        (\Function{LdsExtraTitleMeta}, \Function{LdsExtraHeader}); they
        never reach across the hierarchy to mutate ancestors.
  \item \textbf{One module per document.} Loading two sibling family
        modules at once is unsupported. If a document genuinely sits at
        the intersection of two families (e.g.\ a legal-financial
        addendum), prefer the dominant family's module and add the
        secondary structures inline; do not stack modules to combine
        them.
\end{itemize}

\subsection{Adding a new function class}

\begin{enumerate}
  \item Create \Cmdline{tooling/latex/styles/latex-docs-<family>.sty}.
  \item First non-comment line:
        \Cmdline{\textbackslash RequirePackage\{base.sty\}}.
  \item Register: \Cmdline{\textbackslash LdsRegisterModule\{<key>\}\{<Display Name>\}}.
  \item Add only behaviour shared by ALL documents in this family.
  \item Use \Function{NewLdsCounter} and \Function{NewLdsLabeledBlock}
        for counters and labeled blocks.
  \item Add an example to \Cmdline{src/architecture/style-system/examples/}.
  \item Verify: \Cmdline{cd src/.../examples \&\& latexmk -pdf example-<family>.tex}.
\end{enumerate}

\subsection{Adding a new technical subtype}

Identical to adding a family, except step 2 is
\Cmdline{\textbackslash RequirePackage\{latex-docs-technical\}}, and the
subtype inherits the technical substrate (requirements, hazards, step
results, semantic tagging) without re-declaring it.

\subsection{Backward compatibility}

The canonical house style \Module{latex-docs-style.sty} is unchanged; the
roughly two-hundred existing documents that load it directly continue to
work. The new hierarchy is additive. The parent module loads the canonical
style; family modules load the parent. Concrete documents see one
\Function{usepackage} line, no different in effort from before.

\end{document}



